\subsection*{MT-M35. Finite Euler products become dense after cutoff;
prime exhaustion and cutoff do not commute}

MT-M34 proves actual cutoff-image equality for one Euler
factor.  The following different argument proves equality
of \emph{closed} cutoff images for every nonzero finite
translation polynomial, including Euler products with
incommensurate logarithms.  It does not assert actual-image
surjectivity for these general products.

\paragraph{Finite-translation statement.}
Let $s\ge0$, and set
\[
 H_s=L^2(\mathbf R,(1+x^2)^{s/2}\,dx),\qquad
 U_tf(x)=f(x-t).
 \tag{M35.1}
\]
Let $P$ be a nonzero finite translation polynomial.  Combine
coincident shifts first, delete zero coefficients, and write
\[
                P=\sum_{j=0}^{m}c_jU_{t_j},\qquad
          t_0<t_1<\cdots<t_m,\quad c_j\ne0.
 \tag{M35.2}
\]
Every term is bounded on $H_s$ and $H_{-s}$, by the
polynomial-weight translation bound (M25.7).
Let $\rho\in L^\infty(\mathbf R)$ have essential support
contained in some half-line, in either orientation.  Then
\[
 \overline{M_\rho P H_s}^{\,H_s}
                 =\overline{M_\rho H_s}^{\,H_s}.
 \tag{M35.3}
\]
If $\Theta\subset H_s$ is any dense linear subspace,
the left side is also $\overline{M_\rho P\Theta}$.
For an indicator $\rho=\mathbf1_\Omega$, with $\Omega$
essentially contained in a half-line, this says
\[
 \overline{M_\Omega P\Theta}^{\,H_s}
      =\overline{M_\Omega\overline{P H_s}}^{\,H_s}
      =M_\Omega H_s.
 \tag{M35.4}
\]
In particular it applies to both source cutoffs
$\Omega_{+,L}=[-L,\infty)$ and $\Omega_{0,L}=[-L,L]$.

\paragraph{Endpoint propagation for the weighted dual.}
We first prove the elementary unique-continuation fact
needed here.  Suppose a measurable locally integrable
function $g$ satisfies, almost everywhere,
\[
                 \sum_{j=0}^{m}d_jg(x-\tau_j)=0,
 \quad \tau_0<\cdots<\tau_m,\quad d_j\ne0,
 \tag{M35.5}
\]
and is supported in a half-line.  Then $g=0$ almost
everywhere.  The assertion for $m=0$ is immediate.
For $m\ge1$, translate the equation so that $\tau_0=0$,
and write $H=\tau_m>0$.

If $g=0$ almost everywhere on $(-\infty,A)$, let
$\varepsilon=\tau_1>0$.  For almost every
$x\in[A,A+\varepsilon)$ all the terms with $j\ge1$
evaluate $g$ below $A$.  Since $d_0\ne0$, equation
(M35.5) gives $g(x)=0$ there.  If by induction $g=0$
below $A+n\varepsilon$, the same equation on
$[A+n\varepsilon,A+(n+1)\varepsilon)$ has every
$x-\tau_j<A+n\varepsilon$ for $j\ge1$, and gives
zero on the next interval.  Countably many such steps
prove $g=0$ on the whole line.

If $g=0$ almost everywhere on $(B,\infty)$, put
$\varepsilon'=H-\tau_{m-1}>0$ and evaluate (M35.5)
at $x=y+H$.  It becomes
\[
 d_mg(y)=-\sum_{j<m}d_jg(y+H-\tau_j).
 \tag{M35.6}
\]
For almost every $y\in(B-\varepsilon',B]$, all
arguments on the right exceed $B$.  Induction in the
opposite direction, on consecutive intervals of length
$\varepsilon'$, proves the result.  Finite translations
of the null exceptional sets remain null at each step,
and their countable union is still null.  Combining the
coincident shifts before selecting the two endpoint gaps
is essential to the proof as stated.

We now use the continuous \emph{bilinear} pairing
\[
            \langle f,g\rangle=\int_{\mathbf R}f(x)g(x)\,dx
             \quad (f\in H_s,\ g\in H_{-s}).
 \tag{M35.7}
\]
It identifies $H_{-s}$ with the continuous linear dual
of $H_s$, with complex conjugation absorbed in the
representing function.  Consequently
\[
              P^{\mathrm t}=\sum_{j=0}^{m}c_jU_{-t_j},
              \qquad M_\rho^{\mathrm t}=M_\rho.
 \tag{M35.8}
\]
There is no coefficient conjugation in this convention.
All functions in $H_{-s}$ are locally integrable, since
the weight is bounded above and below on compact sets.

If $g\in H_{-s}$ annihilates $M_\rho PH_s$, then
$P^{\mathrm t}(\rho g)=0$ in $H_{-s}$, hence almost
everywhere.  The function $\rho g$ lies in $H_{-s}$ and
has essential support in a half-line.  Applying
(M35.5)--(M35.6), with the reflected and sorted shifts
from (M35.8), gives $\rho g=0$.  Thus $g$ annihilates
$M_\rho H_s$.  Conversely every annihilator of the
latter image plainly annihilates the former.  Equality
of continuous annihilators gives equality of the closed
linear subspaces, proving (M35.3).  This last assertion
also follows directly from Hilbert orthogonal
complements after converting the bilinear dual to the
usual conjugate-linear Riesz representation.

Boundedness of $M_\rho P$ and density of $\Theta$ prove
the dense-domain assertion.  For an indicator multiplier,
$M_\Omega$ is an orthogonal projection, so its image is
closed.  Boundedness again gives
$\overline{M_\Omega PH_s}=
\overline{M_\Omega\overline{PH_s}}$.
This proves (M35.4).

\paragraph{Finite Euler products and the cofinal cutoff intersection.}
For a finite set $T$ of rational primes, define
\[
                P_T=\prod_{p\in T}(1-U_{-\log p}),
                   \qquad P_\varnothing=1.
 \tag{M35.9}
\]
These are nonzero translation polynomials.  More generally
any product $\prod_v(1-a_vU_{h_v})$, where $h_v>0$ and
$a_v\ne0$, has nonzero constant and largest-shift
coefficients; combining coincident interior shifts cannot
cancel those endpoints.  Reflecting all shifts gives the
same conclusion for (M35.9).  No rational dependence or
independence of the shifts is used in (M35.3).

Let $\Theta$ be the actual trivial compact-character
theta image for $\mathbf Q$, and take $s>1$.
MT-M24--MT-M25, as used in MT-M34, give
$\overline\Theta^{\,H_s}=H_s$ unconditionally from the
classical zero-free boundary of the completed zeta function.
Put
\[
 V_T=\overline{P_T\Theta}^{\,H_s}
                 =\overline{P_TH_s}^{\,H_s}.
 \tag{M35.10}
\]
If $T\subset T'$, the factorization
$P_{T'}=P_TP_{T'\setminus T}$ gives $V_{T'}\subset V_T$.
For either source cutoff and every finite $T$, equation
(M35.4) gives
\[
                  \overline{M_\Omega P_T\Theta}^{\,H_s}
                                  =M_\Omega H_s.
 \tag{M35.11}
\]
Thus for any increasing cofinal sequence $T_k$ of finite
sets of rational primes,
\[
 \bigcap_{k\ge0}
       \overline{M_\Omega P_{T_k}\Theta}^{\,H_s}
                                   =M_\Omega H_s.
 \tag{M35.12}
\]
Every finite stage is the displayed full cutoff space;
there is no exchange of an image with an intersection
in this argument.  Any quotient between two of these
closed cutoff stages is zero.

\paragraph{The uncut cofinal intersection is zero.}
In contrast, for the same $s>1$,
\[
                         \bigcap_{k\ge0}V_{T_k}=\{0\}.
 \tag{M35.13}
\]
Indeed $\int_{\mathbf R}(1+x^2)^{-s/2}dx<\infty$,
so Cauchy--Schwarz gives a continuous embedding
$H_s\subset L^1(\mathbf R)$.  Each functional
\[
                  J_t(f)=\int_{\mathbf R}f(x)e^{itx}\,dx
 \tag{M35.14}
\]
is continuous on $H_s$, and $t\mapsto J_t(f)$ is
continuous by dominated convergence.  The norm-translation
convention gives
\[
       J_t(U_{-\log p}f)=e^{-it\log p}J_t(f).
 \tag{M35.15}
\]
Hence if $p\in T$, every $f\in V_T$ satisfies
$J_{2\pi n/\log p}(f)=0$ for all $n\in\mathbf Z$.
If $f$ belongs to the intersection in (M35.13), cofinality
makes these vanishings valid for every rational prime $p$.

The union of these frequency grids is dense in
$\mathbf R$.  For any fixed $t$, choose primes
$p_j\to\infty$ and choose an integer $n_j$ nearest to
$t\log p_j/(2\pi)$; then
\[
       \left|t-\frac{2\pi n_j}{\log p_j}\right|
                                  \le\frac\pi{\log p_j}\longrightarrow0.
 \tag{M35.16}
\]
Continuity therefore implies $J_t(f)=0$ for every real
$t$.  To recall the uniqueness argument in this setting,
Fourier inversion for a smooth compactly supported test
$\phi$ and Fubini give $\int f\phi=0$: the Fourier
transform of $\phi$ is integrable and $f\in L^1$.
Thus the distribution defined by $f$ vanishes, and
$f=0$ almost everywhere.  This proves (M35.13).

Equations (M35.12)--(M35.13) locate the precise failure
of a cutoff/prime-exhaustion interchange.  For every
fixed finite $L$, the cofinal cutoff intersection is
the full cutoff space; these spaces increase with $L$
and their union is dense in $H_s$.  The uncut cofinal
intersection is zero.  The result concerns subspaces
and their specified closures; it does not identify
the limits of any separately defined kernel operators.

\paragraph{Comparison with the literal Schwartz image in
Theorem 11.1.2.1(3).}
The source's original PDF p.~480 defines the cutoff
spaces in (11.1.2.3) exactly as the closures in (M35.11).
Its cofinal set-theoretic limit (11.1.2.4) is therefore
(M35.12), for $F=\mathbf Q$, $S=\{\infty\}$ and the
trivial character.  Its assertion (3), however, is an
identification with the closure of the specific Schwartz
image (11.1.2.7).  That separate assertion is \emph{not}
refuted by the present example: its stated right side is
also the full cutoff space, as we now check directly.

Every idele class of $\mathbf Q$ has a representative
with positive real component $t=|g|$ and all finite
components in $\mathbf Z_p^\times$: multiply by a rational
number to cancel the finitely many nonzero valuations,
and choose its sign to make the real component positive.
For $S=\{\infty\}$, the sum in (11.1.2.7), before
its cutoff, is then
\[
 \sum_{\xi\in\mathbf Q^\times}
   \left(f_\infty\otimes\prod_p\mathbf1_{\mathbf Z_p^\times}\right)
                   (\xi g)
                          =f_\infty(t)+f_\infty(-t).
 \tag{M35.17}
\]
Indeed only $\xi=1$ and $\xi=-1$ are units at every
finite prime.  This function is in the trivial compact
sector.  The source conditions in this case are
\[
 f_\infty\in\mathcal S(\mathbf R),\qquad
                  f_\infty(0)=0,\qquad
                  \widehat f_\infty(0)=\int_{\mathbf R}f_\infty=0.
 \tag{M35.18}
\]

Fix $L>0$ and take $u\in C_c^\infty(\Omega^\circ)$
for either $\Omega=\Omega_{+,L}$ or $\Omega_{0,L}$.
For $t\ne0$ define
\[
                    f_0(t)=\tfrac12u(\log|t|),
                    \qquad f_0(0)=0.
 \tag{M35.19}
\]
This is an even smooth compactly supported function
whose support is separated from zero, so
$f_0(e^x)+f_0(-e^x)=u(x)$ and $f_0(0)=0$.
Choose an even smooth compactly supported function $r$
with support in $0<|t|<e^{-L}$ and with $\int r\ne0$.
Then
\[
                    f_\infty=f_0-
                            \frac{\int f_0}{\int r}\,r
 \tag{M35.20}
\]
satisfies both conditions in (M35.18).  The repair term
in (M35.17) is supported in $x<-L$, outside either
cutoff.  Consequently the source's cutoff Schwartz
image contains $u$.  Smooth compactly supported
functions in $\Omega^\circ$ are dense in $M_\Omega H_s$:
truncate away from the endpoints and infinity, and then
use ordinary smooth approximation where the weight is
bounded above and below.  Thus the closure of the
right side in (11.1.2.7) is precisely $M_\Omega H_s$,
as claimed here.  No further refutation of part (3)
follows from this rational example.

\paragraph{Input and placement.}
The endpoint-propagation and cutoff-density theorems
are unconditional Hilbert-space deductions, with no
individual-zero or global-Riesz-basis assumption.
The actual theta-density input is the already proved
MT-M24--MT-M25.  Original PDF pp.~428, 480 and 482
were visually checked in MT-M34; the displayed formula
(11.1.2.7) and all its source hypotheses are retained
in the explicit comparison above.  MT-M35 strengthens
the corrected cutoff analysis without declaring the
conditional Hecke-GRH endpoint $M_H$ false or replacing
the remaining geometric-kernel comparison by a bounded
equivariant map.
